\section{Introduzione}
La misurazione accurata dei flussi migratori è essenziale per la definizione di politiche pubbliche basate sulla precisione dei dati.
In Italia, il conteggio degli emigrati si basa sulle cancellazioni anagrafiche per trasferimento di residenza all'estero e sull'iscrizione all'AIRE (Anagrafe degli Italiani Residenti all'Estero)~\cite{legge470_1988}.
Tuttavia, questo sistema soffre di un cronico ritardo nelle registrazioni e di uno scarso senso di adempimento da parte dei cittadini, poiché l'iscrizione non comporta incentivi immediati e le sanzioni previste sono raramente applicate.

Questo fenomeno crea un \textit{esodo invisibile} di individui che, pur risiedendo stabilmente e regolarmente all'estero, risultano ancora formalmente residenti nel territorio nazionale~\cite{migrantes2022, impicciatore_strozza2016}.

L'obiettivo di questo studio è quantificare la discrepanza tra i dati utilizzando la metodologia delle \emph{statistiche speculari} \cite{abel_cohen2019}.
Confrontando quanto dichiarato dall'Italia (ISTAT) con quanto registrato dai paesi di destinazione (Eurostat), identifichiamo non solo la differenza sistematica, ma analizziamo anche le anomalie nella distribuzione temporale delle partenze.